\documentclass[conference]{IEEEtran}
\IEEEoverridecommandlockouts
\usepackage{cite}
\usepackage{amsmath,amssymb,amsfonts}
\usepackage{graphicx}
\usepackage{textcomp}
\usepackage{xcolor}
\usepackage{hyperref}
\usepackage{booktabs}

\begin{document}

\title{Attention Is All You Need}

\author{
\IEEEauthorblockN{Ashish Vaswani\textsuperscript{1}, Noam Shazeer\textsuperscript{1}, Niki Parmar\textsuperscript{2}, Jakob Uszkoreit\textsuperscript{2}, Llion Jones\textsuperscript{2}, Aidan N. Gomez\textsuperscript{3}, Łukasz Kaiser\textsuperscript{1}, Illia Polosukhin\textsuperscript{4}}
\IEEEauthorblockA{\textsuperscript{1}Google Brain, \textsuperscript{2}Google Research, \textsuperscript{3}University of Toronto, \textsuperscript{4}illia.polosukhin@gmail.com}
}

\maketitle

\begin{abstract}
The dominant sequence transduction models are based on complex recurrent or convolutional neural networks that include an encoder and a decoder. The best performing models also connect the encoder and decoder through an attention mechanism. We propose a new simple network architecture, the Transformer, based solely on attention mechanisms, dispensing with recurrence and convolutions entirely. Experiments on two machine translation tasks show these models to be superior in quality while being more parallelizable and requiring significantly less time to train.
\end{abstract}

\begin{IEEEkeywords}
machine learning, sequence transduction, attention mechanisms, transformer
\end{IEEEkeywords}

\section{INTRODUCTION}
Recurrent neural networks, long short-term memory and gated recurrent neural networks in particular, have been firmly established as state of the art approaches in sequence modeling and transduction problems such as language modeling and machine translation. Numerous efforts have since continued to push the boundaries of recurrent language models and encoder-decoder architectures. Recurrent models typically factor computation along the symbol positions of the input and output sequences. Aligning the positions to steps in computation time, they generate a sequence of hidden states as a function of the previous hidden state and the input for position t. This inherently sequential nature precludes parallelization within training examples, which becomes critical at longer sequence lengths, as memory constraints limit batching across examples.

\section{BACKGROUND}
The goal of reducing sequential computation also forms the foundation of the Extended Neural GPU, ByteNet and ConvS2S, all of which use convolutional neural networks as basic building block, computing hidden representations in parallel for all input and output positions. In these models, the number of operations required to relate signals from two arbitrary input or output positions grows in the distance between positions, linearly for ConvS2S and logarithmically for ByteNet. This makes it more difficult to learn dependencies between distant positions. In the Transformer this is reduced to a constant number of operations, albeit at the cost of reduced effective resolution due to averaging attention-weighted positions, an effect we counteract with Multi-Head Attention.

\section{MODEL ARCHITECTURE}
Most competitive neural sequence transduction models have an encoder-decoder structure. Here, the encoder maps an input sequence of symbol representations to a sequence of continuous representations. Given this, the decoder then generates an output sequence of symbols one element at a time. At each step the model is auto-regressive, consuming the previously generated symbols as additional input when generating the next. The Transformer follows this overall architecture using stacked self-attention and point-wise, fully connected layers for both the encoder and decoder.

\subsection{Encoder and Decoder Stacks}
The encoder is composed of a stack of identical layers. Each layer has two sub-layers. The first is a multi-head self-attention mechanism, and the second is a simple, positionwise fully connected feed-forward network. We employ a residual connection around each of the two sub-layers, followed by layer normalization. The decoder is also composed of a stack of identical layers. In addition to the two sub-layers in each encoder layer, the decoder inserts a third sub-layer, which performs multi-head attention over the output of the encoder stack.

\subsection{Attention}
An attention function can be described as mapping a query and a set of key-value pairs to an output, where the query, keys, values, and output are all vectors. The output is computed as a weighted sum of the values, where the weight assigned to each value is computed by a compatibility function of the query with the corresponding key. We call our particular attention "Scaled Dot-Product Attention". The input consists of queries and keys of dimension $d_k$, and values of dimension $d_v$. We compute the dot products of the query with all keys, divide each by $\sqrt{d_k}$, and apply a softmax function to obtain the weights on the values.

\section{WHY SELF-ATTENTION}
In this section we compare various aspects of self-attention layers to the recurrent and convolutional layers commonly used for mapping one variable-length sequence of symbol representations to another sequence of equal length. Motivating our use of self-attention we consider three desiderata. One is the total computational complexity per layer. Another is the amount of computation that can be parallelized, as measured by the minimum number of sequential operations required. The third is the path length between long-range dependencies in the network.

\section{TRAINING}
This section describes the training regime for our models. We trained on the standard WMT 2014 English-German dataset consisting of about 4.5 million sentence pairs. Sentences were encoded using byte-pair encoding, which has a shared source-target vocabulary of about 37000 tokens. For English-French, we used the significantly larger WMT 2014 English-French dataset consisting of 36M sentences and split tokens into a 32000 word-piece vocabulary.

\section{RESULTS}
On the WMT 2014 English-to-German translation task, the big transformer model outperforms the best previously reported models by more than 2.0 BLEU, establishing a new state-of-the-art BLEU score of 28.4. The configuration of this model is listed in the bottom line of Table 3. Training took 3.5 days on 8 P100 GPUs. Even our base model surpasses all previously published models and ensembles, at a fraction of the training cost of any of the competitive models.

\section{MODEL VARIATIONS}
To evaluate the importance of different components of the Transformer, we varied our base model in different ways, measuring the change in performance on English-to-German translation on the development set.

\end{document}